%% 05_methodology.tex — Section 3: Methodology
%% Body-only file — no \documentclass, no \begin{document}.
%% Included via \input{05_methodology} from overleaf/00_main.tex.
%% ─────────────────────────────────────────────────────────────────────────────

\section{Methodology}
\label{sec:methods}

This study addresses the binary classification problem of predicting, prior to
the patient entering the operating room, whether the performed surgical
procedure will deviate from the pre-scheduled procedure.  The proposed pipeline
integrates structured pre-operative features, contextual text representations
derived from bidirectional encoder representations from transformers (BERT)
language models, and a retrieval-augmented mechanism that exploits the
historical deviation behaviour of semantically similar cases.  All feature transformations are estimated
exclusively from training data within each cross-validation fold, thereby
preventing any form of data leakage from validation to training partitions.
The overall workflow is summarised in \cref{fig:flowchart}.

%% ── 3.1 ──────────────────────────────────────────────────────────────────────
\subsection{Data Cleaning and Label Derivation}
\label{sec:cleaning}

The raw dataset ($N = 194{,}661$ surgical cases) was filtered to retain only
cases with a valid pre-operative scheduling timestamp.  All intraoperative
columns (procedure duration, anaesthesia induction and emergence times,
operative diagnosis, and the actually performed procedure) were excluded from
every feature matrix to prevent leakage of intraoperative information into the
pre-operative prediction task.

The binary deviation label was derived as:
\begin{equation}
  y_i = \mathbf{1}\bigl[
        p_i \neq s_i
        \bigr]
  \label{eq:label}
\end{equation}
where $p_i$ denotes the procedure actually performed for case $i$, $s_i$
denotes the procedure scheduled pre-operatively, $y_i = 1$ indicates a
deviation, and $y_i = 0$ indicates no deviation.  After cleaning,
$N = 193{,}995$ cases remained, of which $13{,}134$ ($6.77\,\%$) carried a
positive label, confirming a strongly imbalanced classification problem.

%% ── 3.2 ──────────────────────────────────────────────────────────────────────
\subsection{Structured Feature Engineering}
\label{sec:struct}

The structured pre-operative feature set encompasses four groups.  Demographic
attributes include age at discharge, sex, and mean BMI; continuous variables
were imputed with fold-wise training-set medians to avoid target leakage into
validation partitions.  Clinical severity attributes comprise the ASA physical
status score (ordinal, 1--5) and anaesthetic type, the latter encoded as a
one-hot vector.  Scheduling context attributes include surgical service,
operating room location, and OR trip sequence within the day, all one-hot
encoded.  Temporal attributes include month of year, day of week, and scheduled
start hour, retained as integer features; cyclic encoding was deemed
unnecessary given the dominance of tree-based models.  First and last scheduled
case of day flags, primary procedure status, and surgical encounter type were
included as binary indicators.  Together, these fields yielded 55--56
structured features depending on the fold-specific one-hot cardinality.

%% ── 3.3 ──────────────────────────────────────────────────────────────────────
\subsection{BERT Text Embeddings}
\label{sec:bert}

The scheduled procedure name constitutes the highest-information pre-operative
field: certain procedure types (e.g., exploratory laparoscopy, punch biopsy,
evaluation under anaesthesia) are systematically more likely to expand into a
different definitive procedure than originally booked.  Two pre-trained
transformer models were employed to represent each procedure name as a dense
vector.

ClinicalBERT (Bio-ClinicalBERT), a BERT-base model fine-tuned on MIMIC-III
clinical notes \citep{alsentzer2019publicly}, generated 768-dimensional
\texttt{[CLS]} token embeddings using a batch size of 32 and a maximum
sequence length of 64 tokens.  SentenceBERT (all-MiniLM-L6-v2), a sentence transformer optimised for semantic
similarity through contrastive training \citep{reimers2019sentence}, produced
384-dimensional mean-pooled embeddings using a batch size of 64.  Both embedding
matrices were pre-computed once and cached on disk.

Within each fold, principal component analysis (PCA) was applied to the
training-set embeddings to retain 128 principal components; the same projection
was then applied to the validation set.  Let $\mathbf{E} \in
\mathbb{R}^{n_{\mathrm{train}} \times d}$ denote the training embedding matrix
and $\mathbf{W}_{\mathrm{PCA}} \in \mathbb{R}^{d \times 128}$ the projection
derived from the training partition:
\begin{equation}
  \tilde{\mathbf{e}}_i = \mathbf{W}_{\mathrm{PCA}}^{\top} \mathbf{e}_i
  \label{eq:pca}
\end{equation}
where $\mathbf{e}_i \in \mathbb{R}^{d}$ is the raw embedding of case $i$ and
$\tilde{\mathbf{e}}_i \in \mathbb{R}^{128}$ is its projected representation.
This step reduces dimensionality, accelerates downstream model training, and
mitigates the curse of dimensionality.

%% ── 3.4 ──────────────────────────────────────────────────────────────────────
\subsection{ICD-10 Chapter Encoding}
\label{sec:icd}

The primary pre-operative diagnosis field records the diagnosis in a bracketed
format (e.g., \texttt{(M19.9) arthrosis}).  The first
alphabetic character of each code was extracted and mapped to an
International Classification of Diseases, Tenth Revision (ICD-10) chapter
integer in the range $[1, 22]$ via the standard letter-to-chapter
correspondence (e.g., \texttt{A}/\texttt{B} $\to$ chapter~1, infectious
diseases; \texttt{C}/\texttt{D} $\to$ chapters~2--3, neoplasms; \texttt{M}
$\to$ chapter~13, musculoskeletal disorders).  Cases with an unrecognised or
absent diagnosis code were assigned chapter~0.

This single ordinal feature captures broad diagnostic complexity while avoiding
the high cardinality of full-code one-hot encoding, which would exceed
$10{,}000$ unique values in the present dataset and introduce severe
dimensionality at no demonstrable benefit.

%% ── 3.5 ──────────────────────────────────────────────────────────────────────
\subsection{RAG Augmentation}
\label{sec:rag}

Motivated by retrieval-augmented generation (RAG) \citep{lewis2020retrieval},
the feature vector of each case was enriched with scalar summaries of the
historical deviation behaviour of semantically similar scheduled procedures.
This augmentation provides the classifier with an explicit, data-driven prior
before any model parameters are consulted.

Within each fold, an exact inner-product Facebook AI Similarity Search (FAISS)
index \citep{johnson2019billion} was constructed over the BERT embeddings of
all training cases.  No validation-set embeddings were included in the
index, preventing leakage.  For each case $i$, the $k = 5$ nearest neighbours
$\mathcal{N}(i)$ were retrieved by cosine similarity in embedding space.

Two scalar augmentation features were appended to the feature vector:
\begin{align}
  \phi_{\mathrm{nn}}(i)  &= \frac{1}{k}
      \sum_{j \in \mathcal{N}(i)} y_j
      \label{eq:phi_nn} \\[4pt]
  \phi_{\mathrm{svc}}(i) &=
      \frac{\displaystyle\sum_{j:\,\mathrm{svc}(j)=\mathrm{svc}(i)} y_j}
           {\displaystyle\bigl|\{j : \mathrm{svc}(j) = \mathrm{svc}(i)\}\bigr|}
      \label{eq:phi_svc}
\end{align}
where $\phi_{\mathrm{nn}}(i)$ is the mean deviation rate among the $k$ nearest
training neighbours of case $i$, $\phi_{\mathrm{svc}}(i)$ is the historical
deviation rate of the surgical service assigned to case $i$, $y_j \in \{0,1\}$
is the deviation label of training case $j$, and $\mathrm{svc}(\cdot)$ denotes
the surgical service.

%% ── 3.6 ──────────────────────────────────────────────────────────────────────
\subsection{Fold-wise Encoding and Scaling}
\label{sec:folds}

All feature construction operations, comprising PCA projection, FAISS index
construction and retrieval, fold-wise imputation, and one-hot encoding, were
performed within each fold's training partition.  The resulting transformation
parameters were applied to the corresponding validation partition to preserve
the integrity of the cross-validation protocol.

Six evaluation folds were defined.  Five stratified folds (folds~0--4) each
assigned $80\,\%$ of cases to training ($n \approx 155{,}196$) and $20\,\%$ to
validation ($n \approx 38{,}799$), preserving the class ratio in both
partitions.  A sixth temporal fold (fold~5) used the earliest $80\,\%$ of
cases by scheduling date as the training set and the most recent $20\,\%$ as
the test set; this fold evaluates generalisation to future surgical schedules.
Five feature configurations were evaluated for each fold: structured features
only, ClinicalBERT embeddings, SentenceBERT embeddings, RAG with ClinicalBERT,
and RAG with SentenceBERT.  All features were normalised to $[0, 1]$ using
min-max normalisation fitted on the training partition only.  Encoded feature matrices were serialised to a SQLite database
for reproducible retrieval.

%% ── 3.7 ──────────────────────────────────────────────────────────────────────
\subsection{Model Training and Hyperparameter Optimisation}
\label{sec:models}

Five classifiers were trained and compared across all feature configurations
and folds.  For each model, Optuna \citep{akiba2019optuna} with a
Tree-structured Parzen Estimator (TPE) sampler and a median pruner was used
to conduct 20 hyperparameter trials on a $25\,\%$ random subset of the
training fold.  The best hyperparameters were then used to refit the model on
the complete training partition.

Ridge logistic regression employed an L2 penalty with the L-BFGS solver and
balanced class weighting, with the regularisation strength
$C \in [10^{-3}, 10]$ as the sole tuned parameter.  Random Forest used
balanced class weighting; the number of estimators, maximum tree depth,
maximum features, and minimum samples per split and per leaf were tuned.
XGBoost was executed with GPU acceleration and a fold-dynamic class imbalance
weight $w = n_{-}/n_{+}$, where $n_{-}$ and $n_{+}$ are the counts of
negative and positive training cases respectively; the learning rate, tree
depth, subsampling ratio, column subsampling ratio, and L1/L2 regularisation
terms were tuned.  LightGBM used automatic imbalance correction with early
stopping after 20 rounds of no validation improvement, with analogous
hyperparameters tuned.  The multilayer perceptron (MLP) used the
AdamW optimiser, batch normalisation and dropout after each hidden layer, and a
class-weighted binary cross-entropy loss; the number of hidden layers (1--3),
units per layer, dropout rates, learning rate, and weight decay were tuned,
with training limited to 200 epochs and early stopping with patience 30 on
validation AUC-ROC.

%% ── 3.8 ──────────────────────────────────────────────────────────────────────
\subsection{Performance Evaluation}
\label{sec:eval}

Model performance was assessed using four metrics suited to the strongly
imbalanced class distribution.  The area under the receiver operating
characteristic curve (AUC-ROC) served as the primary threshold-independent
discrimination metric.  The area under the precision-recall curve (AUC-PR),
equivalent to average precision, is more informative than AUC-ROC for
imbalanced problems; a random classifier achieves AUC-PR $\approx 0.068$ on
this dataset.  The F1 score, precision, and recall were computed at a
clinically motivated decision threshold of $\hat{p} = 0.20$, above which a
case is flagged for pre-operative review.  The Brier score was computed to
quantify probability calibration quality.

At the clinical threshold $\hat{p} = 0.20$, a flagged case triggers a set of
pre-operative actions: surgeon briefing, equipment pre-check, extended operating
room block allocation, and anaesthesia team notification.  Performance metrics
are reported as mean $\pm$ standard deviation across folds~0--4, with fold~5
reported separately to assess out-of-time generalisation.
